\chapter{Candidates, physical permission and group selection}\label{ch:menus}
The candidate menu is where an apparently innocent implementation detail
can become an undeclared controller. If the generator proposes only actions
that point toward reference, the later selector is not testing unranked
physical opportunities, however random its final choice may be.

\section{EBU-blind does not mean physically ignorant}
The proposed generator may use the physical information permitted by its
contract: topology, available stock and declared capacity constraints, for
example. It must not rank or shape opportunities using EBU, global potential,
future recovery or an undeclared service objective. Exactly which local
physical information it may read remains a declared design choice.

A finite menu restricts the experiment. If physical quantities are continuous
but the menu contains only a few values, the study concerns that menu,
not all possible divisible action. This is acceptable when explicit. It
becomes misleading when a failure to reach reference is generalized without
checking whether the menu ever made the necessary quantity available.

\section{Joint constraints come before value}
The physical layer validates a complete group. It checks nonnegative
endpoints, conservation, aggregate source and destination constraints,
allowed edges and declared timing. It uses the same rules for the EBU arm
and the physical-random control.

Two individually valid children may exceed a shared cap. If a resolver
changes their quantities, the new accepted vector is the one to value.
After selection, execution must not run another hidden resolver and keep
the old quote. That would settle one vector while performing another.

The order ``physical permission before value'' does not mean physical
permission chooses the most valuable allocation. A proportional allocator,
priority allocator and refusal rule are different policies. A source-locked
study must say which one it uses or avoid unresolved contention in its
declared domain.

\section{Uniform over what?}
Suppose a menu contains physical options A and B, but A appears under two
identifiers. Uniform sampling over entries gives A probability two-thirds;
uniform sampling over unique physical options gives one-half. Neither
interpretation can be inferred from the word random alone.

Group size matters too. Uniform sampling of groups can give an action that
belongs to many groups a high chance of execution. Uniform sampling of
owners first produces another distribution. A no-action option can affect
refusal and time occupation even when it has zero field value. All of these
choices change the transition kernel and must be fixed prospectively.

\section{Local calculation and global assembly}
An exact group value can be evaluated from its affected factors, while the
group assembler still reads a world-wide menu. Thus factor-local valuation
does not prove a decentralized actor process. The specification must name
the chooser, its view, its communication and its authority over owners.

There is no need to settle every distributed-system research question in
the first model. It is necessary to state what was actually chosen. A
centralized experimental sampler can be studied honestly if it is not
advertised as a proof of fully local institutional implementation.

\section*{A useful negative control}
Consider two baseline states with identical permitted local observations
but different hidden remote potential. The candidate generator should not
change its output merely because the hidden diagnostic changed. Designing
such an information-boundary fixture is useful; executing a full transition
with it belongs to the separately authorized integration stage.

\section*{Chapter recap}
The menu, duplicate policy, group assembly and sampling unit are part of
the science. They should not be left to incidental list construction or
silently optimized in the name of convenience.
